% !PRINTED 544
\weberpage{544}
%<<CONT
mon divisor, and consequently, by 3., one can determine the numbers
$\tau_1, \tau_2, \ldots$ in $\mathfrak{o}$ so that
\weq{8}{\varepsilon_1\tau_1 + \varepsilon_2\tau_2 + \cdots = 1}
holds. But from (7), if one multiplies by $\tau_1, \tau_2, \ldots$
and adds, and then sets
\[
\xi = \tau_1\xi_1 + \tau_2\xi_2 + \cdots, \quad
\eta = \tau_1\eta_1 + \tau_2\eta_2 + \cdots
\]
there follows, by means of (8), the equation (5) which was to be proved.

\begin{center}
\S.\,150.

\gesperrt{Fermat's Theorem.}
\end{center}

From the theory of congruences, consequences can be drawn which
correspond exactly to the theorems of rational number theory derived
from \gesperrt{Fermat's theorem}; the most important of these must be
discussed here, for the sake of later application.

By $\pi$ we shall understand a prime functional of degree $f$; that
is, if $p$ is the natural prime divisible by $\pi$, we set
\weq{1}{N_a(\pi) = p^f}
(\S.\,140). If now $\alpha$ is any number in $\mathfrak{o}$ not
divisible by $\pi$, then, as we have already seen in the preceding
section, the product $\alpha\xi$ will run through a complete system of
residues with respect to the modulus $\pi$ at the same time as the
number $\xi$. If we leave out the number divisible by $\pi$, there
remain $N_a(\pi) - 1$ numbers, and if we form the product, it follows
that
\weq{2}{\alpha^{p^f-1}\Pi(\xi) \equiv \Pi(\xi) \ (\mathrm{mod}\ \pi),}
where $\Pi(\xi)$ denotes the product of all the numbers of a complete
system of residues (with the exclusion of zero), and is therefore not
divisible by $\pi$. Accordingly, from (2) follows
\weq{3}{\alpha^{p^f-1} \equiv 1 \ (\mathrm{mod}\ \pi),}
or, if one multiplies by $\alpha$,
\weq{4}{\alpha^{p^f} \equiv \alpha \ (\mathrm{mod}\ \pi),}
and in the last form the theorem holds even when $\alpha$ is divisible
by $\pi$.

The formulae (3), (4), which correspond exactly to
\gesperrt{Fermat's theorem} (Bd.\,I, \S.\,136), shall here too be
called \gesperrt{Fermat's theorem}.
